We will use the following lemma, which characterizes the partition function \eqref{eq:partition_function} in terms of the solution to a Cauchy problem involving the heat equation with initial data $J \in \gmRn$, to prove parts (i) and (ii)(a)-(d) of Prop.~\ref{thm:visc_hj_eqns}.

\begin{lem}[The heat equation with initial data in $\gmRn$] \label{lem:heat_eqn}
Suppose the function $J\colon \Rn \to \R\cup\{+\infty\}$ satisfies assumptions (A1)-(A3).
\begin{enumerate}
\item[(i)] For every $\epsilon>0$, the function $w_\epsilon\colon\Rn\times[0,+\infty)\to (0,1]$ defined by
\begin{equation}
w_{\epsilon}(\bx,t)\coloneqq \frac{1}{(2\pi t \epsilon)^{n/2}}Z_J(\bx,t,\epsilon) = \frac{1}{\left(2\pi t\epsilon\right)^{n/2}}\int_{\Rn}e^{-\left(\frac{1}{2t}\left\Vert \bx-\by\right\Vert _{2}^{2}+J(\by)\right)/\epsilon}\diff\by\label{eq:heat_w}
\end{equation}
is the unique smooth solution to the Cauchy problem
\begin{equation}
\begin{dcases}
\frac{\partial w_{\epsilon}}{\partial t}(\bx,t)=\frac{\epsilon}{2}\Laplacian w_{\epsilon}(\bx,t) & \mbox{in }\Rn\times(0,+\infty),\\
w_{\epsilon}(\bx,0)=e^{-J(\bx)/\epsilon} & \mbox{in }\Rn.
\end{dcases}\label{eq:cauchy_heat_prob}
\end{equation}
In addition, the domain of integration of the integral \eqref{eq:heat_w} can be taken to be $\dom J$ or, up to a set of Lebesgue measure zero, $\inter(\dom J)$
or $\dom\partial J$. Furthermore, for every $\bx\in\Rn$ and $\epsilon > 0$, except possibly at the points $\bx\in(\dom J)\setminus(\inter\dom J)$ if such points exist, the pointwise limit of $w_\epsilon(\bx,t)$ as $t \to 0$ exists and satisfies
\begin{equation*}
\lim_{\substack{t\to0\\ t>0}}w_{\epsilon}(\bx,t)=e^{-J(\bx)/\epsilon},
\end{equation*}
with the limit equal to $0$ whenever $\bx\notin\dom J$. If $\bx\in(\dom J)\setminus\inter(\dom J)$, then we may only conclude 
\[
\limsup_{\substack{t\to0\\t>0}}w_{\epsilon}(\bx,t)\leqslant e^{-J(\bx)/\epsilon}
\]
and
\[
\liminf_{\substack{t\to0\\t>0}}w_{\epsilon}(\bx,t)\geqslant e^{-J(\bx)/\epsilon}\left(\frac{1}{(2\pi\epsilon)^{n/2}}\int_{\dom J}e^{-\frac{1}{2\epsilon}\left\Vert \bx-\by\right\Vert _{2}^{2}}\diff\by\right).
\]

\item[(ii)]  (Log-concavity and monotonicity properties).
\begin{enumerate}
\item The function $\Rn\times(0,+\infty)\ni(\bx,t)\mapsto t^{n/2}w_{\epsilon}(\bx,t)$ is jointly log-concave.
\item The function $(0,+\infty)\ni t\mapsto t^{n/2}w_{\epsilon}(\bx,t)$ is strictly monotone increasing.
\item The function $(0,+\infty)\ni\epsilon\mapsto\epsilon^{n/2}w_{\epsilon}(\bx,t)$ is strictly monotone increasing.
\item The function $\Rn\ni\bx\mapsto e^{\frac{1}{2t\epsilon}\normsq{\bx}}w_{\epsilon}(\bx,t)$ is strictly log-convex.
\end{enumerate}
\end{enumerate}
\end{lem}
The proof of (i) follows from classical PDEs arguments for the Cauchy problem \eqref{eq:cauchy_heat_prob} tailored to the initial data $(\bx,\epsilon)\mapsto e^{-J(\bx)/\epsilon}$ with $J$ satisfying assumptions (A1)-(A3), and the proof of log-concavity and monotonicity (ii)(a)-(d) follows from the Pr\'ekopa--Leindler and H\"older's inequalities \citep{leindler1972certain, prekopa1971logarithmic,folland2013real}; we present the details below.


\begin{proof}
\textbf{Proof of Lemma~\ref{lem:heat_eqn} (i):} For every $\bx \in \Rn$, $t>0$, and $\epsilon>0$, let $w_\epsilon(\bx,t) = \frac{1}{\left(2\pi t\epsilon\right)^{n/2}} \int_{\dom J} e^{-\left(\frac{1}{2t}\left\Vert \bx-\by\right\Vert _{2}^{2}+J(\by)\right)/\epsilon}\diff\by$. By assumptions (A1)-(A3) that $J$ is convex with $\inter(\dom J) \neq \varnothing$ and bounded from below by zero, we have
\[
0 \leqslant \frac{1}{\left(2\pi t\epsilon\right)^{n/2}} \int_{\dom J} e^{-\left(\frac{1}{2t}\left\Vert \bx-\by\right\Vert _{2}^{2}+J(\by)\right)/\epsilon}\diff\by \leqslant \frac{1}{\left(2\pi t\epsilon\right)^{n/2}} \int_{\dom J} e^{-\frac{1}{2t\epsilon}\left\Vert \bx-\by\right\Vert _{2}^{2}}\diff\by = 1,
\]
i.e., $w_\epsilon(\bx,t)$ is finite, non-negative and bounded from above by one. If $\dom J \neq \Rn$, the domain of integration can still be taken to be $\Rn$ with the understanding that $e^{-J(\by)/\epsilon} = 0$ for every $\by\notin\dom J$. The domain of integration can also be taken to be $\inter(\dom J)$ or $\dom\partial J$ since the boundary points of the domain of $J$ is a set of Lebesgue measure zero relative to $\Rn$ (see definition \ref{def:convex_sets}). Now, a routine application of the theorem of differentiation under the integral sign~\cite{folland2013real} implies that the function $w_\epsilon\colon\Rn\times[0,+\infty) \to (0,1]$ is smooth, and a short calculation shows that it solves the heat equation for every $\bx \in \Rn$ and $t > 0$. Uniqueness of solutions follows by non-negativity of $w_\epsilon(\bx,t)$ for every $\bx \in \Rn$, $t>0$ and $\epsilon = 0$ \cite{donnelly1987uniqueness}.

We will now use Fatou's lemma (\citep{folland2013real}, Lemma 2.18) to compute bounds for the two limits $\limsup_{\substack{t\to0\\t>0}}w_{\epsilon}(\bx,t)$ and $\liminf_{\substack{t\to0\\t>0}}w_{\epsilon}(\bx,t)$ for every $\bx\in\Rn$ and $\epsilon>0$. First, note that the change of variables $\frac{\bx-\by}{\sqrt{t}}\mapsto\by$ in \eqref{eq:heat_w} yields
\[
w_{\epsilon}(\bx,t)=\frac{1}{(2\pi\epsilon)^{n/2}}\int_{\Rn}e^{-\left(\frac{1}{2}\left\Vert \by\right\Vert _{2}^{2}+J(\bx-\sqrt{t}\by)\right)/\epsilon}\diff\by,
\]
so that the function
\[
\by\mapsto e^{-\frac{1}{2\epsilon}\left\Vert y\right\Vert _{2}^{2}}\left(1-e^{-J(\bx-\sqrt{t}\by)/\epsilon}\right)
\]
is non-negative. The reverse Fatou's lemma therefore applies to this function, and hence
\[
\limsup_{\substack{t\to0\\t>0}}w_{\epsilon}(\bx,t)\leqslant\frac{1}{(2\pi\epsilon)^{n/2}}\int_{\Rn}e^{-\frac{1}{2\epsilon}\left\Vert \by\right\Vert _{2}^{2}}\left(\limsup_{\substack{t\to0\\t>0}}e^{-J(\bx-\sqrt{t}\by)/\epsilon}\right)\diff\by.
\]
Using the lower semicontinuity of $J$, the limit inside the integral satisfies
\begin{eqnarray*}
\limsup_{\substack{t\to0\\t>0}}e^{-J(\bx-\sqrt{t}\by)/\epsilon} & = & e^{\limsup_{\substack{t\to0\\t>0}}-J(\bx-\sqrt{t}\by)/\epsilon}\\
& = & e^{-\liminf_{\substack{t\to0\\t>0}}J(\bx-\sqrt{t}\by)/\epsilon}\\
& \leqslant & e^{-J(\bx)/\epsilon},
\end{eqnarray*}
and therefore $\limsup_{\substack{t\to0\\t>0}}w_{\epsilon}(\bx,t)\leqslant e^{-J(\bx)/\epsilon}$ for every $\bx\in\Rn$. If $\bx\notin\dom J$, then
\[
0\leqslant\liminf_{\substack{t\to0\\t>0}}e^{-J(\bx-\sqrt{t}\by)/\epsilon}\leqslant\limsup_{\substack{t\to0\\t>0}}e^{-J(\bx-\sqrt{t}\by)/\epsilon}\leqslant0,
\]
which implies $\lim_{\substack{t\to0\\t>0}}w_{\epsilon}(\bx,t)=0$ for every $\bx\notin\dom J$. Suppose now that $\bx\in\dom J$. By Fatou's lemma, 
\[
\liminf_{\substack{t\to0\\t>0}}w_{\epsilon}(\bx,t)\geqslant\frac{1}{(2\pi\epsilon)^{n/2}}\int_{\Rn}e^{-\frac{1}{2\epsilon}\left\Vert \by\right\Vert _{2}^{2}}\left(\liminf_{\substack{t\to0\\t>0}}e^{-J(\bx-\sqrt{t}\by)/\epsilon}\right)\diff\by.
\]
If $\bx\in\inter(\dom J)$, then by continuity $\lim_{\substack{t\to0\\t>0}}J(\bx-\sqrt{t}\by)=J(\bx)$ for every $\by\in\Rn$, and so $\liminf_{\substack{t\to0\\t>0}}w_{\epsilon}(\bx,t)\geqslant e^{-J(\bx)/\epsilon}$. Combined with the $\limsup$, we find $\lim_{\substack{t\to0\\t>0}}w_{\epsilon}(\bx,t)=e^{-J(\bx)/\epsilon}$ for every $\bx \in \inter(\dom J)$ and $\bx \notin \dom J$. If $\bx\in(\dom J)\setminus\inter(\dom J)$, if any such point exists, and $0<t\leqslant1$, then convexity of $J$ implies
\begin{align*}
J(\bx-\sqrt{t}\by) & =J((1-\sqrt{t})\bx+\sqrt{t}(\bx-\by))\\
& \leqslant(1-\sqrt{t})J(\bx)+\sqrt{t}J(\bx-\by),
\end{align*}
and hence for any $0<t\leqslant1$,
\begin{align*}
w_{\epsilon}(\bx,t) & \geqslant e^{-(1-\sqrt{t})J(\bx)/\epsilon}\frac{1}{(2\pi\epsilon)^{n/2}}\int_{\Rn}e^{-\frac{1}{2\epsilon}\left\Vert \by\right\Vert _{2}^{2}}e^{-\sqrt{t}J(\bx-\by)/\epsilon}\diff\by,\\
& =e^{-(1-\sqrt{t})J(\bx)/\epsilon}\frac{1}{(2\pi\epsilon)^{n/2}}\int_{\Rn}e^{-\frac{1}{2\epsilon}\left\Vert \bx-\by\right\Vert _{2}^{2}}e^{-\sqrt{t}J(\by)/\epsilon}\diff\by.
\end{align*}
Since $\liminf_{\substack{t\to0\\t>0}}e^{-\sqrt{t}J(\by)/\epsilon} = 1$ for every $\by\in\dom J$ and is equal to $0$ for every $\by\notin\dom J$, Fatou's lemma yields
\[
\liminf_{\substack{t\to0\\t>0}}w_{\epsilon}(\bx,t)\geqslant e^{-J(\bx)/\epsilon}\left(\frac{1}{(2\pi\epsilon)^{n/2}}\int_{\dom J}e^{-\frac{1}{2\epsilon}\left\Vert \bx-\by\right\Vert _{2}^{2}}\diff\by\right)
\]
for every $\bx\in(\dom J)\setminus\inter(\dom J)$, if any such point exists.

\textbf{Proof of Lemma~\ref{lem:heat_eqn} (ii)(a):} The log-concavity property will be shown using the Pr\'ekopa--Leindler inequality.
\begin{thm}\label{thm:prekopa}
[Pr\'ekopa--Leindler inequality \citep{leindler1972certain, prekopa1971logarithmic}]
Let $f$, $g,$ and $h$ be non-negative real-valued and measurable functions on $\Rn$, and suppose
\[
h(\lambda\by_{1}+(1-\lambda)\by_{2})\geqslant f(\by_{1})^{\lambda}g(\by_{2})^{(1-\lambda)}
\]
for every $\by_{1},$ $\by_{2}\in\Rn$ and $\lambda\in(0,1)$. Then
\[
\int_{\Rn}h(\by)\diff\by\geqslant\left(\int_{\Rn}f(\by)\diff\by\right)^{\lambda}\left(\int_{\Rn}g(\by)\diff\by\right)^{(1-\lambda)}.
\]
\end{thm}
\textbf{Proof of Lemma~\ref{lem:heat_eqn} (ii)(a) (continued):} Let $\epsilon>0$, $\lambda\in(0,1)$, $\bx=\lambda\bx_{1}+(1-\lambda)\bx_{2}$, $\by=\lambda\by_{1}+(1-\lambda)\by_{2}$, and $t=\lambda t_{1}+(1-\lambda)t_{2}$ for any $\bx_{1},\,\bx_{2},\,\by_{1},\,\by_{2}\in\Rn$ and $t_{1},\,t_{2}\in(0,+\infty)$. The joint convexity of the function $\Rn\times(0,+\infty)\ni(\boldsymbol{z},t)\mapsto\frac{1}{2t}\left\Vert \boldsymbol{z}\right\Vert _{2}^{2}$ and convexity of $J$ imply
\[
\frac{1}{2t}\left\Vert \bx-\by\right\Vert _{2}^{2}+J(\by) \leqslant \frac{\lambda}{2t_{1}}\left\Vert \bx_{1}-\by_{1}\right\Vert _{2}^{2}+\frac{(1-\lambda)}{2t_{2}}\left\Vert \bx_{2}-\by_{2}\right\Vert _{2}^{2}+\lambda J(\by_{1})+(1-\lambda)J(\by_{2}),
\]
This gives
\[
\frac{e^{-\left(\frac{1}{2t}\left\Vert \bx-\by\right\Vert _{2}^{2}+J(\by)\right)/\epsilon}}{(2\pi\epsilon)^{n/2}}\geqslant\left(\frac{e^{-\left(\frac{1}{2t_{1}}\left\Vert \bx_{1}-\by_{1}\right\Vert _{2}^{2}+J(\by_{1})\right)/\epsilon}}{(2\pi\epsilon)^{n/2}}\right)^{\lambda}\left(\frac{e^{-\left(\frac{1}{2t_{2}}\left\Vert \bx_{2}-\by_{2}\right\Vert _{2}^{2}+J(\by_{2})\right)/\epsilon}}{(2\pi\epsilon)^{n/2}}\right)^{1-\lambda}.
\]
Applying the Pr\'ekopa--Leindler inequality with
\[
h(\by)=\frac{e^{-\left(\frac{1}{2t}\left\Vert \bx-\by\right\Vert _{2}^{2}+J(\by)\right)/\epsilon}}{(2\pi\epsilon)^{n/2}},
\]
\[
f(\by)=\frac{e^{-\left(\frac{1}{2t_{1}}\left\Vert \bx_{1}-\by\right\Vert _{2}^{2}+J(\by)\right)/\epsilon}}{(2\pi\epsilon)^{n/2}},
\]
and 
\[
g(\by)=\frac{e^{-\left(\frac{1}{2t_{2}}\left\Vert \bx_{2}-\by\right\Vert _{2}^{2}+J(\by)\right)/\epsilon}}{(2\pi\epsilon)^{n/2}},
\]
and using the definition \eqref{eq:heat_w} of $w_{\epsilon}(\bx,t)$, we get
\[
t^{n/2}w_{\epsilon}(\bx,t)\geqslant\left(t_{1}^{n/2}w_{\epsilon}(\bx_{1},t_{1})\right)^{\lambda}\left(t_{2}^{n/2}w_{\epsilon}(\bx_{2},t_{2})\right)^{(1-\lambda)},
\]
As a result, the function $(\bx,t)\mapsto t^{n/2}w_{\epsilon}(\bx,t)$ is jointly log-concave on $\Rn\times(0,+\infty)$.

\textbf{Proof of Lemma~\ref{lem:heat_eqn} (ii)(b):} Since $t\mapsto\frac{1}{t}$ is strictly monotone decreasing on $(0,+\infty)$, then for every $\bx\in\Rn$, $\by \in \dom J$, $\epsilon>0$, and $0<t_1<t_2$,
\[
\frac{e^{-\left(\frac{1}{2t_1}\left\Vert \bx-\by\right\Vert _{2}^{2}+J(\by)\right)/\epsilon}}{(2\pi\epsilon)^{n/2}}<\frac{e^{-\left(\frac{1}{2t_2}\left\Vert \bx-\by\right\Vert _{2}^{2}+J(\by)\right)/\epsilon}}{(2\pi\epsilon)^{n/2}}
\]
whenever $\bx\neq\by$. Integrating both sides of the inequality with respect to $\by$ over $\dom J$ yields
\[
\frac{1}{(2\pi\epsilon)^{n/2}}\int_{\dom J}e^{-\left(\frac{1}{2t_1}\left\Vert \bx-\by\right\Vert _{2}^{2}+J(\by)\right)/\epsilon}\diff\by<\frac{1}{(2\pi\epsilon)^{n/2}}\int_{\dom J}e^{-\left(\frac{1}{2t_2}\left\Vert \bx-\by\right\Vert _{2}^{2}+J(\by)\right)/\epsilon}\diff\by,
\]
As a result, the function $t\mapsto t^{n/2}w_{\epsilon}(\bx,t)$ is strictly monotone increasing on $(0,+\infty)$. 

\textbf{Proof of Lemma~\ref{lem:heat_eqn} (ii)(c):} Since $\epsilon\mapsto\frac{1}{\epsilon}$ is strictly monotone decreasing on $(0,+\infty)$ and $\dom J \ni \by\mapsto J(\by)$ is non-negative by assumption (A3), then for every $\bx\in\Rn$, $t>0$, and $0<\epsilon_1<\epsilon_2$ we have
\[
e^{-\left(\frac{1}{2t}\left\Vert \bx-\by\right\Vert _{2}^{2}+J(\by)\right)/\epsilon_1}<e^{-\left(\frac{1}{2t}\left\Vert \bx-\by\right\Vert _{2}^{2}+J(\by)\right)/\epsilon_2}
\]
whenever $\bx\neq\by$. Integrating both sides of the inequality with respect to $\by$ over $\dom J$ yields 
\[
\int_{\dom J}e^{-\left(\frac{1}{2t}\left\Vert \bx-\by\right\Vert _{2}^{2}+J(\by)\right)/\epsilon_1}d\by<\int_{\dom J}e^{-\left(\frac{1}{2t}\left\Vert \bx-\by\right\Vert _{2}^{2}+J(\by)\right)/\epsilon_2}\diff\by,
\]
As a result, the function $\epsilon\mapsto\epsilon^{n/2}w_{\epsilon}(\bx,t)$ is strictly monotone increasing on $(0,+\infty)$.

\textbf{Proof of Lemma~\ref{lem:heat_eqn} (ii)(d):} Let $\epsilon>0$, $t>0$, $\lambda\in(0,1)$, $\bx_1,\bx_2\in \Rn$ with $\bx_1\neq\bx_2$ and $\bx=\lambda\bx_{1}+(1-\lambda)\bx_{2}$. Then
\begin{align*}
    e^{\frac{1}{2t\epsilon}\normsq{\bx}}w_\epsilon(\bx,t) &= \frac{1}{(2\pi t\epsilon)^{n/2}}\int_{\dom J} e^{\left(\left\langle \bx,\by\right\rangle/t -\frac{1}{2t}\normsq{\by} - J(\by)\right)/\epsilon}\diff\by \\
    &=\int_{\dom J} \left(\frac{e^{\left(\left\langle \bx_1,\by\right\rangle/t-\frac{1}{2t}\normsq{\by} - J(\by)\right)/\epsilon}}{(2\pi t \epsilon)^{n/2}}\right)^\lambda \left(\frac{e^{\left\langle \bx_2,\by\right\rangle/t\epsilon-\frac{1}{2t}\normsq{\by} - J(\by)/\epsilon}}{(2\pi t \epsilon)^{n/2}}\right)^{1-\lambda} \diff\by.
\end{align*}
H\"older's inequality (\citep{folland2013real}, Theorem 6.2) then implies
\begin{align*}
    e^{\frac{1}{2t\epsilon}\normsq{\bx}}w_\epsilon(\bx,t) &\leqslant \left(\int_{\dom J} \frac{e^{\left(\left\langle \bx_1,\by\right\rangle/t-\frac{1}{2t}\normsq{\by} - J(\by)\right)/\epsilon}}{(2\pi t \epsilon)^{n/2}} \diff\by\right)^\lambda \left(\int_{\dom J}  \frac{e^{\left\langle \bx_2,\by\right\rangle/t\epsilon-\frac{1}{2t}\normsq{\by} - J(\by)/\epsilon}}{(2\pi t \epsilon)^{n/2}} \diff\by \right)^{1-\lambda} \\
    &=\left(e^{\frac{1}{2t\epsilon}\normsq{\bx_1}}w_\epsilon(\bx_1,t)\right)^\lambda \left(e^{\frac{1}{2t\epsilon}\normsq{\bx_2}}w_\epsilon(\bx_2,t)\right)^{1-\lambda},
\end{align*}
where the inequality in the equation above is an equality if and only if there exists a constant $\alpha \in \mathbb{R}$ such that $\alpha e^{\left\langle \bx_1,\by\right\rangle/t\epsilon} = e^{\left\langle \bx_x,\by\right\rangle/t\epsilon}$
for almost every $\by \in \dom J$. This does not hold here since $\bx_1\neq\bx_2$. As a result, the function $\Rn\ni\bx\mapsto e^{\frac{1}{2t\epsilon}\normsq{\bx}}w_{\epsilon}(\bx,t)$ is strictly log-convex.
\end{proof}


\textbf{Proof of Prop.~\ref{thm:visc_hj_eqns} (i) and (ii)(a)-(d):} The proof of these statements follow from Lemma~\ref{lem:heat_eqn} and classic results about the Cole--Hopf transform (see, e.g., \citep{evans1998partial}, Section 4.4.1), with $S_\epsilon(\bx,t) \coloneqq -\epsilon\log(w_\epsilon(\bx,t))$.

\textbf{Proof of Prop.~\ref{thm:visc_hj_eqns} (iii):} The formulas follow from a straightforward calculation of the gradient, divergence, and Laplacian of $S_\epsilon(\bx,t)$ that we omit here. Since the function $\bx \mapsto \frac{1}{2}\normsq{\bx}-tS_{\epsilon}(\bx,t)$ is strictly convex, we can invoke (Corollary 26.3.1, \cite{rockafellar1970convex}) to conclude that its gradient $\bx \mapsto \bx - t\nabla_{\bx}S_\epsilon(\bx,t)$, which gives the posterior mean $\bu_{PM}(\bx,t,\epsilon)$, is bijective.

\textbf{Proof of Prop.~\ref{thm:visc_hj_eqns} (iv):} We will prove this result in three steps. First, we will show that 
\[
\limsup_{\substack{\epsilon\to0\\
\epsilon>0
}
}S_{\epsilon}(\bx,t)\leqslant\inf_{\by\in\inter(\dom J)}\left\{ \frac{1}{2t}\left\Vert \bx-\by\right\Vert _{2}^{2}+J(\by)\right\} 
\]
 and 
\[
\inf_{\by\in\inter(\dom J)}\left\{ \frac{1}{2t}\left\Vert \bx-\by\right\Vert _{2}^{2}+J(\by)\right\} = \inf_{\by\in\dom J}\left\{ \frac{1}{2t}\left\Vert \bx-\by\right\Vert _{2}^{2}+J(\by)\right\} \equiv S_0(\bx,t).
\]
Next, we will show that $\liminf_{\substack{\epsilon\to0\\
\epsilon>0
}
}S_{\epsilon}(\bx,t)\geqslant S_0(\bx,t)$. 
Finally, we will use steps 1 and 2 to conclude that $\lim_{\substack{\epsilon\to0\\
\epsilon>0
}
}S_{\epsilon}(\bx,t)=S_0(\bx,t)$. Pointwise and local uniform convergence of the gradient $\lim_{\substack{\epsilon\to0\\
\epsilon>0
}
}\nabla_{\bx}S_{\epsilon}(\bx,t)=\nabla_{\bx}S_0(\bx,t)$, the partial derivative $\lim_{\substack{\epsilon\to0\\
\epsilon>0
}
}\frac{\partial S_{\epsilon}(\bx,t)}{\partial t}=\frac{\partial S_0(\bx,t)}{\partial t}$, and the Laplacian $\lim_{\substack{\epsilon\to0\\
\epsilon>0}}\frac{\epsilon}{2}\Laplacian S_{\epsilon}(\bx,t)=0$ then follow from the convexity and differentiability of the solutions $(\bx,t)\mapsto S_0(\bx,t)$ and $(\bx,t)\mapsto S_{\epsilon}(\bx,t)$ to the HJ PDEs~\eqref{eq:hj_non_visc_pde} and~\eqref{eq:hj_visc_pde}.

In what follows, we will use the following large deviation principle result \citep{deuschel2001large}: For every Lebesgue measurable set $\mathcal{A} \in \Rn$,
\[
\lim_{\substack{\epsilon\to0\\
\epsilon>0
}
}-\epsilon\ln\left(\frac{1}{(2\pi t\epsilon)^{n/2}}\int_{\mathcal{A}}e^{-\frac{1}{2t\epsilon}\left\Vert \bx-\by\right\Vert _{2}^{2}}\diff\by\right)=\essinf_{\by\in\mathcal{A}}\frac{1}{2t}\left\Vert \bx-\by\right\Vert _{2}^{2},
\]
where essential infimum means infimum that holds almost everywhere, i.e., if the essential infimum above is attained at $\boldsymbol{a}\in\mathcal{A}$, then $\frac{1}{2t}\left\Vert \bx-\boldsymbol{a}\right\Vert _{2}^{2}\leqslant\frac{1}{2t}\left\Vert \bx-\by\right\Vert _{2}^{2}$ for almost every $\by\in\mathcal{A}$. 

\textbf{Step 1.} (Adapted from Deuschel and Stroock~\cite{deuschel2001large}, Lemma 2.1.7.) By convexity, the function $J$ is continuous for every $\by_{0}\in\inter(\dom J)$, the latter set being open. Therefore, for every such $\by_{0}$ there exists a number $r_{\by_{0}}>0$ such that for every $0<r\leqslant r_{\by_{0}}$ the open ball $B_{r}(\by_{0})$ is contained in $\inter(\dom J)$. Hence
\begin{align*}
S_{\epsilon}(\bx,t) & \coloneqq-\epsilon\ln\left(\frac{1}{(2\pi t\epsilon)^{n/2}}\int_{\inter(\dom J)}e^{-(\frac{1}{2t}\normsq{\bx-\by}+J(\by))/\epsilon}\diff\by\right)\\
 & \leqslant-\epsilon\ln\left(\frac{1}{(2\pi t\epsilon)^{n/2}}\int_{B_{r}(\by_{0})}e^{-(\frac{1}{2t}\normsq{\bx-\by}+J(\by))/\epsilon}\diff\by\right)\\
 & \leqslant-\epsilon\ln\left(\frac{1}{(2\pi t\epsilon)^{n/2}}\int_{B_{r}(\by_{0})}e^{-\frac{1}{2t\epsilon}\normsq{\bx-\by}}\diff\by\right)+\sup_{\by\in B_{r}(\by_{0})}J(\by).
\end{align*}
Take $\limsup_{\substack{\epsilon\to0\\ \epsilon>0}}$ and apply the large deviation principle to the term on the right to get
\[
\limsup_{\substack{\epsilon\to0\\\epsilon>0}}S_{\epsilon}(\bx,t)\leqslant\essinf_{\by\in B_{r}(\by_{0})}\frac{1}{2t}\left\Vert \bx-\by\right\Vert _{2}^{2}+\sup_{\by\in B_{r}(\by_{0})}J(\by).
\]
Take $\lim_{r\to0}$ on both sides of the inequality to find
\[
\limsup_{\substack{\epsilon\to0\\\epsilon>0}}S_{\epsilon}(\bx,t)\leqslant\frac{1}{2t}\left\Vert \bx-\by_{0}\right\Vert _{2}^{2}+J(\by_{0}).
\]
Since the inequality holds for every $\by_{0}\in\inter(\dom J)$, we can take the infimum over all $y\in\inter(\dom J)$ on the right-hand-side of the inequality to get
\begin{equation}\label{eq:limsupineq1}
\limsup_{\substack{\epsilon\to0\\\epsilon>0}}S_{\epsilon}(\bx,t)\leqslant\inf_{\by\in\inter(\dom J)}\left\{ \frac{1}{2t}\left\Vert \bx-\by\right\Vert _{2}^{2}+J(\by)\right\}.
\end{equation}
By assumptions (A1) and (A2) that $J \in \gmRn$ and $\inter(\dom J) \neq \varnothing$, the infimum on the right hand side is equal to that taken over $\dom J$ (\citep{rockafellar1970convex}, Corollary 7.3.2), i.e.,
\begin{equation}\label{eq:limsupineq2}
\inf_{\by\in\inter(\dom J)}\left\{ \frac{1}{2t}\left\Vert \bx-\by\right\Vert _{2}^{2}+J(\by)\right\}   =  \inf_{\by\in\dom J}\left\{ \frac{1}{2t}\left\Vert \bx-\by\right\Vert _{2}^{2}+J(\by)\right\} \equiv S_0(\bx,t).
\end{equation}
We combine~\eqref{eq:limsupineq1} and~\eqref{eq:limsupineq2} to obtain 
\[\limsup_{\substack{\epsilon\to0\\\epsilon>0}}S_{\epsilon}(\bx,t)\leqslant S_0(\bx,t),
\]
which is the desired result.

\textbf{Step 2.} We can invoke Lemma 2.1.8 in~\cite{deuschel2001large} because its conditions are satisfied (in the notation of \cite{deuschel2001large}, $\Phi=-J$, which is upper semicontinuous, $\by\mapsto\frac{1}{2t}\left\Vert \bx-\by\right\Vert _{2}^{2}$ is the rate function, and note that the tail condition (2.1.9) is satisfied in that $\sup_{\by\in\Rn} -J(\by) = -\inf_{\by\in\Rn} J(\by) = 0$ by assumption (A3)) to get
\[
\liminf_{\substack{\epsilon\to0\\\epsilon>0}}S_{\epsilon}(\bx,t)\geqslant S_0(\bx,t).
\]

\textbf{Step 3.} Combining the two limits derived in steps 1 and 2 yield
\[
\lim_{\substack{\epsilon\to0\\\epsilon>0}}S_{\epsilon}(\bx,t)=S_0(\bx,t)
\]
for every $\bx\in\Rn$ and $t>0$, where the limit converges uniformly on every compact subset $(\bx,t)$ of $\Rn\times(0,+\infty)$ (\citep{rockafellar1970convex}, Theorem 10.8). 

% Include domains
By differentiability and joint convexity of both $\Rn \times (0,+\infty) \ni (\bx,t) \mapsto S_0(\bx,t)$ and $\Rn \times (0,+\infty) \ni (\bx,t)\mapsto S_{\epsilon}(\bx,t)-\frac{n\epsilon}{2}\ln t$ (Prop.~\ref{thm:e_u_nonvisc_hj} (i), and Prop.~\ref{thm:visc_hj_eqns} (i) and (ii)(a)), we can invoke (\citep{rockafellar1970convex}, Theorem 25.7) to get 
\[
\lim_{\substack{\epsilon\to0\\\epsilon>0}}\nabla_{\bx}S_{\epsilon}(\bx,t)=\nabla_{\bx}S_0(\bx,t)\,\mbox{and}\,\lim_{\substack{\epsilon\to0\\\epsilon>0}}\left(\frac{\partial S_{\epsilon}(\bx,t)}{\partial t}-\frac{n\epsilon}{2t}\right)=\lim_{\substack{\epsilon\to0\\
\epsilon>0}}\frac{\partial S_{\epsilon}(\bx,t)}{\partial t}=\frac{\partial S_0(\bx,t)}{\partial t},
\]
for every $\bx\in\Rn$ and $t>0$, where the limit converges uniformly on every compact subset of $\Rn\times(0,+\infty)$. Furthermore, the viscous HJ PDE~\eqref{eq:hj_visc_pde} for $S_\epsilon$ implies that
\begin{align*}
\lim_{\substack{\epsilon\to0\\
\epsilon>0
}
}\frac{\epsilon}{2}\Laplacian S_{\epsilon}(\bx,t) & =\lim_{\substack{\epsilon\to0\\
\epsilon>0
}
}\left(\frac{\partial S_{\epsilon}(\bx,t)}{\partial t}+\frac{1}{2}\left\Vert \nabla_{\bx}S_{\epsilon}(\bx,t)\right\Vert ^{2}\right),\\
 & =\left(\frac{\partial S_{0}(\bx,t)}{\partial t}+\frac{1}{2}\left\Vert \nabla_{\bx}S_0(\bx,t)\right\Vert ^{2}\right)\\
 & =0,
\end{align*}
where the last equality holds thanks to the HJ PDE~\eqref{eq:hj_non_visc_pde} (see Prop.~\ref{thm:e_u_nonvisc_hj}). Here, again, the limit holds for every $\bx\in\Rn$ and $t>0$, and the limit converges uniformly over any compact subset of $\Rn\times(0,+\infty)$. Finally, the limit $\lim_{\substack{\epsilon\to0\\\epsilon>0}} \bu_{PM}(\bx,t,\epsilon) = \bu_{MAP}(\bx,t)$ holds directly as a consequence to the limit $\lim_{\substack{\epsilon\to0\\\epsilon>0}}\nabla_{\bx}S_{\epsilon}(\bx,t)=\nabla_{\bx}S_0(\bx,t)$ and the representation formulas~\eqref{eq:grad_s_eps} (see Prop.~\ref{thm:visc_hj_eqns}(iii)) and \eqref{eq:grad_map_rep} (see Prop.~\ref{thm:e_u_nonvisc_hj}(ii))for the posterior mean and MAP estimates, respectively.